\pg{349} II. Dada a hypothese da existencia real de um só ente finito, deduz-se legitimamente a existencia do Ente infinito \footnote{Ente simples é 0 que não tem partes. --- Se se considera a razão ou o fundamento da \emph{simplicidade} do ente, este pode ser, como dizemos no texto, \emph{negativo}, \emph{positivo}, e \emph{abstractivo} Ou \emph{precisivo}. --- Se se considera a qualidade das partes, que se excluem do ente simples, este divide-se principalmente, e em opposição ao \emph{composto}, em \emph{physico}, \emph{metaphaysico} e logico. conforme as parles, que se excluem, são \emph{physicas} (materia e forma), ou \emph{metaphysicas} (potencia e acto. essencia e existencia), ou \emph{logtcas} (genero e differenga). --- Se se considera o grau da \emph{simplicidade} do ente, este, como dizemos no texto, pode ser \emph{absoluto} e \emph{relativo}, conforme exelue --- toda e qualquer especie de \emph{composição} (\textsc{Deus}, é só Elle, é \emph{absolutamente} simples) --- ou somente uma \emph{ou} outra especie (\emph{a} nossa alma é \emph{relativamente} simples, porque, se exclue a composição \emph{physica}, não exclue a \emph{metaphysica}),}.

-. AgrTIGO TII

\emph{Ente simples e ente composto}.

79, Ente simples e ente composto. --- Ente \emph{szmples} é o que não tem partes, e por isso é indivisivel. Pode ser \emph{negativo}, \emph{positivo} e \emph{abstractivo}, conforme pode não ter partes --- ou pela tenuidade e imperfeição da sua entidade, como é o ponto mathematico, --- ou pela perfeição do seu ser, como é a nossa alma, --- ou pelo effeito da abstracgão, como é o ente em geral. Neste artigo occupamo-nos do \emph{simples positivo}. --- Ente \emph{composto} é o que tem partes, e por isso é divisivel. Pode ser \emph{physico}, \emph{metaphysico} e \emph{logico}, conforme as partes, são --- \emph{physicas}, como a alma, e o corpo, --- ou meftaphysicas; como-a potencia e o acto; a essencia e a existencia, --- ou \emph{logicas}; como o genero e a dilferença e

73. O ente simples, quando subsiste perfeitamente, é \textbf{mais} perfeito que o composto. --- O ente \emph{simples}, que subsiste perfeitamente, por não constar de partes, tem todos os elementos ou graus, que constituem a sua entidade, não só unidos, mas indivisiveis; ao passo que o ente \emph{composto}, embora tenha os seus elementos unidos, comtudo não os tem indivisiveis. Ora um ente,

(t) Nem se diga que entre o finito e o infinito n\&o existe proporção, e por isso o conhecimento d'aquelle não pode levar ao conhecimento d'este. Por quanto, se entre o finito e o infinito não existe uma proporção de \emph{entidade} e de egualdade, existe uma proporção de dependencia e de analogia; pois o finito depende do infinito, como o effeito depende da causa. --- Esta falla de proporção de \emph{entidade} e de egualdade obsta a que tenhamos um conhecimento proprio e adequado do infinito; mas não obsta a que conhe- «amos a existencia do mesmo infinito, e, de algum modo, as suas perteigóes.

